\documentclass[10pt,aspectratio=169]{beamer}

\usetheme{metropolis}
\usepackage[utf8]{inputenc}
\usepackage[T1]{fontenc}
\usepackage[english]{babel}
\usepackage{listings}
\usepackage{booktabs}
\usepackage{amsmath,amssymb}

\definecolor{gfblue}{HTML}{1F4E79}
\definecolor{codebg}{HTML}{F4F4F2}
\definecolor{warn}{HTML}{B03A2E}
\definecolor{good}{HTML}{1E6B3A}

\setbeamercolor{frametitle}{bg=gfblue}
\setbeamercolor{alerted text}{fg=warn}

\lstset{
  basicstyle=\ttfamily\scriptsize,
  backgroundcolor=\color{codebg},
  frame=none,
  columns=fullflexible,
  keepspaces=true,
  showstringspaces=false,
  breaklines=true,
  xleftmargin=2mm,
  literate=
    {á}{{\'a}}1 {é}{{\'e}}1 {í}{{\'i}}1 {ý}{{\'y}}1 {ú}{{\'u}}1
    {č}{{\v{c}}}1 {ě}{{\v{e}}}1 {ň}{{\v{n}}}1 {ř}{{\v{r}}}1
    {š}{{\v{s}}}1 {ž}{{\v{z}}}1 {ů}{{\r{u}}}1
    {å}{{\aa}}1 {ä}{{\"a}}1 {ö}{{\"o}}1 {ü}{{\"u}}1
    {è}{{\`e}}1 {ê}{{\^e}}1 {ł}{{\l}}1 {ę}{{\k{e}}}1 {ą}{{\k{a}}}1
    {→}{{$\rightarrow$}}1 {∀}{{$\forall$}}1 {⊢}{{$\vdash$}}1 {∈}{{$\in$}}1
}

\newcommand{\bad}[1]{\textcolor{warn}{#1}}
\newcommand{\ok}[1]{\textcolor{good}{#1}}
\newcommand{\commit}[1]{\texttt{\footnotesize #1}}
\newcommand{\prompt}[1]{\begin{quote}\small\itshape ``#1''\end{quote}}

% the three standard slides of a task
\newcommand{\why}{\par\textcolor{gfblue}{\textbf{Why}}\par}
\newcommand{\how}{\par\textcolor{gfblue}{\textbf{Prompt and interaction}}\par}
\newcommand{\what}{\par\textcolor{gfblue}{\textbf{Outcome}}\par}

\title{Twenty-three Commits}
\subtitle{What Claude Code has contributed to Informath, July--September 2026}
\author{Aarne Ranta}
\date{September 2026}

\begin{document}

\maketitle

\begin{frame}{This talk}
  \begin{itemize}
    \item Informath: 899 commits since 2024, almost all by hand.
    \item \textbf{23} of them are co-authored by Claude (Claude Code, Anthropic's
          coding agent, working in a terminal in my checkout).
    \item For each task: \textbf{why} it was needed, the \textbf{prompt and
          interaction}, and the \textbf{outcome} --- including what is still wrong.
    \item Sources: \texttt{git log}, the diffs, Claude Code session transcripts
          (September), session logs that Claude wrote itself (July).
  \end{itemize}

  \vspace{0.5em}
  \small
  These slides were themselves drafted by Claude from those sources ---
  a 24th task, not committed yet.
\end{frame}

\begin{frame}{Outline}
  \small\tableofcontents
\end{frame}

% =====================================================================
\section{Informath in five minutes}

\begin{frame}{The problem}
  Mathematics exists in two registers.

  \begin{columns}[T]
    \begin{column}{0.48\textwidth}
      \textbf{Formal}
      \begin{itemize}
        \item Agda, Rocq, Lean, Dedukti
        \item machine-checkable
        \item unreadable to most humans
      \end{itemize}
    \end{column}
    \begin{column}{0.48\textwidth}
      \textbf{Informal}
      \begin{itemize}
        \item English, French, German, Swedish, \dots
        \item readable
        \item ambiguous, unverifiable
      \end{itemize}
    \end{column}
  \end{columns}

  \vspace{1em}
  Informath translates between them in all directions:
  \textbf{informalization} (formal $\to$ informal),
  \textbf{autoformalization} (informal $\to$ formal),
  and informal $\to$ informal via the formal.
\end{frame}

\begin{frame}[fragile]{One statement, many languages}
\begin{lstlisting}
Dedukti: prop110 : (a : Elem Int) -> (c : Elem Int) ->
  Proof (and (odd a) (odd c)) ->
  Proof (forall Int (b => even (plus (times a b) (times b c)))).
\end{lstlisting}

\begin{itemize}\footnotesize
  \item \textbf{En}: Let $a$ and $c$ be integers. Assume that $a$ and $c$ are odd.
        Then $ab + bc$ is even for all integers $b$.
  \item \textbf{Fr}: Soient $a$ et $c$ des entiers. Supposons que $a$ et $c$ sont impairs.
        Alors $ab+bc$ est pair pour tous les entiers $b$.
  \item \textbf{Sv}: Låt $a$ och $c$ vara heltal. Anta att $a$ och $c$ är udda.
        Då är $ab + bc$ jämnt för alla heltal $b$.
\end{itemize}

\footnotesize All generated from the Dedukti source, and Agda, Rocq and Lean as well.
\end{frame}

\begin{frame}{Architecture}
  \begin{center}
  \begin{tabular}{c}
    \textbf{Dedukti} --- the semantic core \\
    $\updownarrow$ \\
    \textbf{MathCore} --- abstract syntax close to the logic \\
    $\updownarrow$ \quad (semantics / NLG in Haskell) \\
    \textbf{Informath} --- variation, synonyms, symbolic vs.\ verbal \\
    $\updownarrow$ \quad (GF concrete syntaxes, using the RGL) \\
    \textbf{Eng \quad Fre \quad Ger \quad Swe} \quad (+ Latex, + experimental Fin Cze Pol) \\
  \end{tabular}
  \end{center}

  \vspace{0.5em}
  \begin{itemize}\small
    \item GF: Grammatical Framework; RGL: its Resource Grammar Library (40+ languages).
    \item The grammar is split into a language-independent \emph{functor}
          and small per-language modules of words and function words.
  \end{itemize}
\end{frame}

\begin{frame}[fragile]{Symbol tables: from constants to words}
  Each Dedukti constant gets a wording and/or a notation, given \emph{by example}:
\begin{lstlisting}
isGroup     : "#1 is a group"
isSubgroup  : "#1 is a subgroup of #2"
inv         : "the inverse of #2" | $#2^{-1}$
cosetRel    : "#3 is congruent to #4 modulo #2"
\end{lstlisting}
  \begin{itemize}\small
    \item The example is \emph{parsed} with the grammar into a lexical tree
          (\texttt{NounPrepNoun2 subgroup\_Noun of\_Prep}), which then
          works in any statement, in any language.
    \item \texttt{\$...\$} forms become LaTeX macros.
    \item Only shapes the grammar can parse are allowed --- a recurring theme today.
  \end{itemize}
\end{frame}

\begin{frame}{How Claude worked}
  \begin{itemize}
    \item Claude Code runs in a terminal, reads and edits the files, runs
          \texttt{gf}, \texttt{stack}, \texttt{RunInformath}, \texttt{pdflatex}.
    \item I give short prompts in English; Claude plans, works, reports.
    \item Commits only when asked; I push myself.
    \item It can spawn \emph{subagents} for parallel subtasks (e.g.\ batches of words).
    \item Remembers preferences across sessions in a small memory file,
          e.g.\ ``prefer symbolic expressions''.
  \end{itemize}
\end{frame}

\begin{frame}{The 23 commits as 11 tasks}
  \footnotesize
  \begin{tabular}{rllr}
    \toprule
    & \textbf{Task} & \textbf{Date} & \textbf{Commits} \\
    \midrule
    1 & Finnish words and function words & 10 July & 2 \\
    2 & A Czech grammar & 10, 16 July & 4 \\
    3 & A Polish grammar & 16 July & 2 \\
    4 & Completing the French Wikidata lexicon & 22 July & 1 \\
    \midrule
    5 & Words for the Godement formalization & 24--25 Sept & 5 \\
    6 & Clash-free macro names & 24 Sept & 1 \\
    7 & Proof texts & 25 Sept & 4 \\
    8 & Compressing hypotheses in proof texts & 26 Sept & 1 \\
    9 & Formulas in standard logical notation & 26 Sept & 1 \\
    10 & GodementSyntax: the language of definitions and theorems & 26 Sept & 1 \\
    11 & GodementExamples: richer symbol tables & 26 Sept & 1 \\
    \bottomrule
  \end{tabular}

  \vspace{0.5em}
  Part I: new languages. \quad Part II: formalizing Godement's \emph{Algebra}.
\end{frame}

% =====================================================================
\section{Part I: New languages}

\begin{frame}{Task 1: Finnish --- why}
  \begin{itemize}
    \item A Finnish grammar existed as a \emph{compilable template},
          made by copying Swedish and English.
    \item It said \texttt{linking ... OK} --- and it said
          \texttt{then\_Adv = mkAdv "så"}.
    \item \texttt{VerbalConstantsFin}: an English copy with two Finnish words.
    \item \texttt{WikidataWordsFin} (2000+ terms) was a byte-identical copy
          of English and did not compile, so it was commented out:
  \end{itemize}
  \begin{quote}\small
    a grammar that ``worked'' by excluding 2000 of its 2100 words.
  \end{quote}
  \begin{itemize}
    \item The question: can an LLM complete a GF grammar for a language
          it knows well, given the RGL?
  \end{itemize}
\end{frame}

\begin{frame}{Task 1: Finnish --- prompt and interaction}
  \small Goal given (as recorded in Claude's session log):
  \prompt{Complete the Finnish concrete grammars by replacing placeholder
    (mostly Swedish / English-copy) content with proper Finnish mathematical
    vocabulary.}
  \begin{itemize}
    \item The prompts themselves were not saved; this was the first experiment.
    \item Three layers: $\sim$23 function words, $\sim$120 core math words,
          2051 Wikidata terms.
    \item The Wikidata terms went to \textbf{7 parallel subagents}, each returning
          an English--Finnish table; a Python script generated GF.
    \item Claude found why the lexicon did not compile, and dropped the dependency
          on the Finnish dictionary: Finnish \emph{smart paradigms} suffice.
    \item My interventions: asking about prepositions; not accepting
          \texttt{linking ... OK} as proof of correctness.
  \end{itemize}
\end{frame}

\begin{frame}{Task 1: Finnish --- outcome}
  \commit{9ac0431}, \commit{ab82b63} (10 July)
  \begin{itemize}
    \item 2051 entries: 1216 nouns, 749 adjectives, 86 verbs; all enabled; links.
    \item \ok{\emph{lyhde}} (sheaf), \ok{\emph{ydin}} (kernel), \ok{\emph{jaoton}} (prime);
          \emph{olkoon $x$ ryhmä}; \texttt{mkNoun "rengas"} gives
          \emph{kaikille renkaille $R$} with consonant gradation for free.
  \end{itemize}
  \textbf{Remaining problems}
  \begin{itemize}\small
    \item Five prepositions are loose words (\texttt{mkPrep "yli"}) where Finnish uses cases:
          ``They compile. They read acceptably. They are not right.''
    \item Invented terms: \bad{quiver $\to$ \emph{nuolikko}}; capitalized
          \emph{Abelin} next to \emph{abelin}; 33 orphans, 15 duplicates inherited from English.
    \item No native review yet; moved to \texttt{grammars/next/}.
  \end{itemize}
\end{frame}

\begin{frame}{Task 2: Czech --- why}
  \begin{itemize}
    \item Czech: a new Informath language, and a test of a \emph{weak} RGL.
    \item The template was Swedish down to the morphology:
          \texttt{mkN "typ" "typer"}.
    \item The Czech RGL was a stub: \texttt{SentenceCze} had 4 rules,
          no \texttt{mkV}, no \texttt{V3}; missing rules crashed only late:
          \texttt{NOT YET IMPLEMENTED: ComplVS}.
    \item Informath's functor calls the RGL by qualified names, so the gaps
          could not be worked around inside Informath.
  \end{itemize}
\end{frame}

\begin{frame}{Task 2: Czech --- prompt and interaction}
  \prompt{Make \texttt{InformathCze.gf} compile, link, and produce Czech.}
  \begin{itemize}
    \item A loop: crash, write the missing RGL rule, rebuild the library ---
          about 20 times. ``Tirelessly.''
    \item Half of the session was spent in the RGL (commits there too).
    \item 2003 terms by \textbf{8 subagents}, with French and German as hints.
    \item Asked to audit itself, Claude re-implemented the RGL's paradigm
          selector in Python and compared: \bad{225 of 1177 nouns (19\%)}
          fell back to a wrong declension --- two RGL bugs, affecting every Czech grammar.
    \item A human question: ``is that actually an imperative?''
  \end{itemize}
\end{frame}

\begin{frame}{Task 2: Czech --- outcome}
  \commit{e81159b}, \commit{f756218} (10 July), \commit{114fbd4}, \commit{ab7e423} (16 July)
  \begin{itemize}\small
    \item InformathCze compiles and links; its own verb conjugator for
          \emph{-ovat, -at, -it/-et/-ět}.
    \item 2003 terms, nouns with nominative, genitive and gender; wrong
          fallbacks from 225 to 1 after 7 new paradigms.
    \item 789 words in Godement's register; three-place verbs after an RGL fix.
  \end{itemize}
  \textbf{Remaining problems}
  \begin{itemize}\small
    \item \bad{\emph{existuje matice M, takový, že}} --- should be \emph{taková}:
          ``such that'' cannot agree in gender.
    \item Multi-word nouns decline only the head; ``imperatives'' are present tense.
    \item UserExtensionsCze still disabled; $\sim$1500 terms unreviewed.
  \end{itemize}
\end{frame}

\begin{frame}{Task 3: Polish --- why}
  \begin{itemize}
    \item The second Slavic language, now with Czech as the model.
    \item The Polish RGL was in good shape (only 6 missing rules) ---
          the problem turned out to be \emph{nouns}.
    \item A test of whether the lessons of Czech generalize.
  \end{itemize}
\end{frame}

\begin{frame}{Task 3: Polish --- prompt and interaction}
  \prompt{Create the Polish concrete syntax modules, following the Cze
    structure, so that InformathPol.gf compiles.}
  and then, on request, Mizar words and user extensions.
  \begin{itemize}\small
    \item Three dead ends with nouns:
      \begin{itemize}\footnotesize
        \item \texttt{mkN} guessed \bad{\emph{zbióra, elementowie, przypadeka}};
        \item the two-argument \texttt{mkN} dropped the genitive; fixing it gave
              \bad{\emph{liczbaa}} --- reverted and documented;
        \item the 1057 SGJP tables bake in an example word:
              \emph{samolot} turns \emph{wektor} into \bad{\emph{wektot}}.
      \end{itemize}
    \item So Claude wrote its own declensions, taking the stem from the genitive.
    \item 12 subagents; they \emph{refused} to fake genitives for adjectival nouns,
          which led to a new paradigm.
    \item Random generation found \bad{\emph{dla wszystkich płytki}}; a test of
          \texttt{meet\_Verb2} found a word-order bug in the RGL.
  \end{itemize}
\end{frame}

\begin{frame}{Task 3: Polish --- outcome}
  \commit{6ee105f}, \commit{c735985} (16 July)
  \begin{itemize}\small
    \item InformathPol links: 2003 terms, 0 rejects;
          \ok{\emph{niech $x$ będzie grupą}}.
    \item Mizar words reuse Wikidata terms where the sense matches:
          \emph{pole relacji}, not \emph{ciało}.
    \item The session log ends with six lessons ``for the next language''.
  \end{itemize}
  \textbf{Remaining problems}
  \begin{itemize}\small
    \item \bad{\emph{istnieje macierz M, taki, że}} --- the same agreement problem as Czech.
    \item \bad{\emph{koł}} for \emph{kół}; weak terms (\emph{pushforwardowy}).
    \item Needs a hand-compiled \texttt{MarkupPol}: \texttt{make next\_grammar} fails
          on a fresh machine.
  \end{itemize}
\end{frame}

\begin{frame}{Task 4: French Wikidata --- why and interaction}
  \why
  \begin{itemize}\small
    \item French was an official language, but \texttt{WikidataWordsFre} was
          hand-curated: 537 of 2003 terms, no verbs.
  \end{itemize}
  \how
  \prompt{Complete \texttt{WikidataWordsFre.gf} so that every abstract
    WikidataWords fun has a French linearization.}
  \begin{itemize}\small
    \item A diff: 1481 missing (1008 nouns, 387 adjectives, 86 verbs).
    \item \textbf{15 parallel agents}, $\sim$100 terms each; 15 name collisions dropped.
    \item A second pass on multi-word plurals: \emph{trous spectraux},
          \emph{chemins de fer}.
    \item One smooth session.
  \end{itemize}
\end{frame}

\begin{frame}[fragile]{Task 4: French Wikidata --- outcome}
  \commit{9699638} (22 July): 2003/2003 terms, compiles.
\begin{lstlisting}
lin close_Verb = mkVerb (mkV "fermer") ; ---Claude
\end{lstlisting}
  \begin{itemize}\small
    \item Every new line is marked \texttt{-{}-{}-Claude} (1466 of them):
          \texttt{grep -v} gives a lexicon without them until they are checked.
  \end{itemize}
  \textbf{Remaining problems}
  \begin{itemize}\small
    \item Uncertain senses: \emph{stack} $\to$ \emph{champ}, \emph{net} $\to$ \emph{filet},
          \bad{\emph{quiver} $\to$ \emph{carquois}} (a quiver for arrows).
    \item $\sim$21 invariant multi-word adjectives; source typos carried over
          (\texttt{minmum\_Noun}).
    \item All unreviewed. Next candidates: German ($\sim$1650 missing),
          Swedish ($\sim$1850).
  \end{itemize}
\end{frame}

% =====================================================================
\section{Part II: Godement's \emph{Algebra}}

\begin{frame}{The Godement project}
  \begin{itemize}
    \item Roger Godement, \emph{Algebra} (1968): a classic textbook,
          Part I (sets) and Part II (groups, rings, modules, linear algebra).
    \item Since 18 September, a companion repository \texttt{godement-index}:
      \begin{itemize}\small
        \item the whole book as a UD treebank: 10\,413 sentences;
        \item a Dedukti formalization of Part I and §§ 6--24, with Claude,
              each constant linked to Godement's own sentence;
        \item symbol tables, English readings, formulas, and an HTML index.
      \end{itemize}
    \item The Informath commits of Part II were \emph{needed by} this work:
          missing words, bugs found, new outputs, new grammar.
  \end{itemize}
\end{frame}

\begin{frame}{Task 5: Words for Godement --- why and interaction}
  \why
  \begin{itemize}\small
    \item Every symbol-table entry must parse. Godement's wordings needed words
          the lexicon lacked: ``$R$ is \emph{undecidable}'', ``$X$ is
          \emph{equipotent} to $Y$'', ``a \emph{two-sided} ideal'',
          ``a \emph{Cramer} system''.
  \end{itemize}
  \how
  \prompt{let us continue from PARTI-FORMALIZATION.md to complete the Godement
    book till the end of chapter 5}
  \begin{itemize}\small
    \item Claude added a word whenever a table entry failed to parse, and
          committed it with the formalization (\emph{``make a commit if ready for it; I will push''}).
    \item It deliberately did \emph{not} add \texttt{countable\_Adj}: already in
          WikidataWords, and a second declaration breaks the build.
  \end{itemize}
\end{frame}

\begin{frame}[fragile]{Task 5: Words for Godement --- outcome}
  \commit{583e0e5}, \commit{a4427ce} (24 Sept), \commit{bc4d431}, \commit{00c01cb},
  \commit{6509853} (25 Sept)
\begin{lstlisting}
fun 'two-sided_Adj' : Adj ;
lin 'two-sided_Adj' = mkAdj "two-sided" ;
\end{lstlisting}
  \begin{itemize}\small
    \item 9 words: \emph{undecidable, contradictory, equipotent, generate,
          two-sided, finite-dimensional, alternating, increasing, Cramer}.
    \item Symbol tables for Part I and §§ 6--24 build.
  \end{itemize}
  \textbf{Remaining problems}
  \begin{itemize}\small
    \item Words were the easy part: the \emph{shapes} of entries were limited
          --- no three-place nouns, no ``into'', no negation (see Task 11).
  \end{itemize}
\end{frame}

\begin{frame}{Task 6: Clash-free macro names --- why}
  \begin{itemize}
    \item A LaTeX form \texttt{\$...\$} in a symbol table becomes a macro.
    \item Its name was the letters of the constant + \texttt{MACRo}:
          digits and underscores were dropped.
    \item \texttt{card0}, \texttt{card1}, \texttt{card2} all became
          \texttt{\textbackslash cardMACRo}; the last definition silently won.
    \item Nothing warned. The document compiled and said
      \begin{center}
        \bad{$2 \cdot x = 2$} \quad for \quad $0 \cdot x = 0$, \qquad
        \bad{$x^2 = x$} \quad for \quad $x^1 = x$.
      \end{center}
  \end{itemize}
\end{frame}

\begin{frame}{Task 6: Clash-free macro names --- prompt and interaction}
  \begin{itemize}\small
    \item Found while reading Part I back. First: \emph{``what are you doing?''}
          and \emph{``I see that it stalls. Can you just jump over that?''}
    \item Claude proposed two fixes: spell digits as words
          (\texttt{\textbackslash cardZeroMACRo}), or rename the constants.
  \end{itemize}
  \prompt{I know these fixes. They would not solve everything, because there
    can be names that already contain Zero etc. It would be better to generate
    all macro names with an algorithm that guarantees uniqueness. For instance,
    suffixes MACRoA, MACRoB, \dots{} MACRoAA, MACRoAB, \dots}
  \begin{itemize}\small
    \item Implemented as a bijective base-26 index over the whole table;
          \emph{``this sounds good''}.
    \item Claude corrected its own diagnosis twice after measuring
          (the stall was one oversized theorem, not the flags).
  \end{itemize}
\end{frame}

\begin{frame}{Task 6: Clash-free macro names --- outcome}
  \commit{f445edd} (24 Sept)
  \begin{center}\small
  \begin{tabular}{lll}
    \toprule
    & \textbf{before} & \textbf{after} \\
    \midrule
    macros & \texttt{\textbackslash cardMACRo} (3 constants)
           & \texttt{\textbackslash cardMACRoY, ...Z, ...AA} \\
    \texttt{mulCard\_zero} & \bad{$2 \cdot x = 2$} & \ok{$0 \cdot x = 0$} \\
    \texttt{powCard\_one} & \bad{$x^2 = x$} & \ok{$x^1 = x$} \\
    \texttt{projections} & \bad{$\mathrm{pr}_2(x,y) = x$} & \ok{$\mathrm{pr}_1(x,y) = x$} \\
    \bottomrule
  \end{tabular}
  \end{center}
  \textbf{Remaining problems}
  \begin{itemize}\small
    \item Names depend on position: adding a LaTeX form renames later macros,
          so derived files must be rebuilt together.
    \item Informath has no regression test suite; checking is by comparing outputs.
  \end{itemize}
\end{frame}

\begin{frame}{Task 7: Proof texts --- why}
  \begin{itemize}
    \item Godement's theorems now had real Dedukti \emph{proofs} (§ 6 onwards).
    \item Informath could print a proof as numbered lines, but:
      \begin{itemize}\small
        \item every argument of a rule became a line, mostly
              ``there is a proposition'': § 6 Theorem 1 gave 46 lines, 14 of them steps;
        \item the output was a LaTeX 2.09 demo (\texttt{\textbackslash documentstyle})
              that did not compile;
        \item lines were raw token lists: \texttt{\$ y \$ is equal to \$ x \$ .}
      \end{itemize}
    \item Goal: proofs that a mathematician can read, as PDF.
  \end{itemize}
\end{frame}

\begin{frame}{Task 7: Proof texts --- prompt and interaction}
  \begin{itemize}
    \item No transcript survives for this day; the story is in four commits
          made within half an hour:
      \begin{enumerate}\small
        \item show the steps, not the arguments (a heuristic filter);
        \item replace the heuristic by \emph{profiles}:
              \texttt{\#SHOW forallE [3,4]} in the symbol table, or read the
              premisses off the rule's type;
        \item a document that can be read and compiled;
        \item unlex the lines.
      \end{enumerate}
    \item Each was committed together with the regenerated proofs in
          \texttt{godement-index}.
    \item A slip: commit 2 swept in unrelated files (slides, logs,
          \texttt{*.gf\textasciitilde} backups).
  \end{itemize}
\end{frame}

\begin{frame}[fragile]{Task 7: Proof texts --- outcome}
  \commit{0623e2b}, \commit{dabdd50}, \commit{d183036}, \commit{2ef5f55} (25 Sept)
\begin{lstlisting}
5. Assume that x is equal to y.
6. Assume that y is equal to z.
7. x is equal to y and y is equal to z.        andI 5, 6
8. x is equal to z.                            ifE 4, 7
\end{lstlisting}
  \begin{itemize}\small
    \item PDFs of the proofs for §§ 6--10, 12--17, 22 (27 pages for § 6).
  \end{itemize}
  \textbf{Remaining problems}
  \begin{itemize}\small
    \item Each proof is headed by its raw Dedukti type, not an English theorem.
    \item \bad{§ 6 Theorem 1 prints two lines}: the declaration, then its
          \emph{conclusion} as an assumption; the steps are lost.
    \item Proof labels leak into formulas (``the value of $f$ at \texttt{hy}'').
  \end{itemize}
\end{frame}

\begin{frame}{Task 8: Compressing hypotheses --- why and interaction}
  \why
  \begin{itemize}\small
    \item Every variable had its own line: ``Let $x$ be \dots{} Let $y$ be \dots{}
          Let $z$ be \dots'' --- long and unlike mathematical prose.
  \end{itemize}
  \how
  \prompt{I want to improve src/ProofText.hs by compressing hypotheses. If
    there are consecutive proof lines introducing a variable with the same type,
    they should be compressed to one line using the abstract syntax function
    GVarsHypo on the list of those variables. The lines should of course be
    renumbered to reflect this.}
  \begin{itemize}\small
    \item Claude also found that the Dedukti table (\texttt{-proof-terms})
          would no longer match, and that the flag had never been accepted at all.
    \item \emph{``Do compress the Dedukti table. But where in the code are the wrong
          variable names created?''} --- diagnosed (beta reduction under the rule's
          binder), fix proposed, not made.
  \end{itemize}
\end{frame}

\begin{frame}[fragile]{Task 8: Compressing hypotheses --- outcome}
  \commit{c391de6} (26 Sept)
  \begin{columns}[T]
    \begin{column}{0.47\textwidth}
      \textbf{Before} (8 lines)
\begin{lstlisting}
1. Let x be a mathematical object.
2. Let y be a mathematical object.
3. Let z be a mathematical object.
4. If x is equal to y and ...
\end{lstlisting}
    \end{column}
    \begin{column}{0.47\textwidth}
      \textbf{After} (6 lines)
\begin{lstlisting}
1. Let x, y and z be mathematical
   objects.
2. If x = y and y = z, then x = z.
3. Assume that x = y.
\end{lstlisting}
    \end{column}
  \end{columns}
  \begin{itemize}\small
    \item The symbolic $x = y$ came from another change; asked about it:
          \emph{``Yes, I did want to increase the usage of symbolic expressions
          instead of verbal ones.''} --- saved as a lasting preference.
  \end{itemize}
  \textbf{Remaining problems}
  \begin{itemize}\small
    \item Wrong bound-variable names (``for all elements $x$ and $m$ of $\mathbb{N}$'')
          are still there.
  \end{itemize}
\end{frame}

\begin{frame}{Task 9: Standard logical notation --- why and interaction}
  \why
  \begin{itemize}\small
    \item Readers want each statement also as a compact formula
          ($\forall$, $\land$, $\to$, $\vdash$), next to the English.
    \item Informath had symbolic notations for constants, but not for the logic.
  \end{itemize}
  \how
  \prompt{I would like to have grammars/MathCoreLatex.gf and InformathLatex that
    generates formulas with standard logical notation and uses the symbolic
    notations defined in SymbolicConstantsLatex and in .dkgf files. In
    RunInformath, the flag -to-symbolic-latex should translate Dedukti to this
    notation. [\dots] Use the precedences in the same way as in TermsLatex.}
  \begin{itemize}\small
    \item Claude chose conventions itself and offered to change them: sequents
          $\Gamma \vdash A$, $:\Leftrightarrow$ for definitions, no proofs shown.
    \item \emph{``Also update godement-index.html with these formulas.''}
  \end{itemize}
\end{frame}

\begin{frame}{Task 9: Standard logical notation --- outcome}
  \commit{f87c660} (26 Sept): three new GF modules, a flag, a Makefile target.
  \begin{itemize}\small
    \item English: ``Let $f$ be a function from mathematical objects to mathematical
          objects. Let $u$ and $v$ be mathematical objects. Assume that $u = v$.
          Then $f(u) = f(v)$.''
    \item Formula: \quad $f \in \mathrm{Obj} \to \mathrm{Obj},\ u, v \in \mathrm{Obj},\ u = v
          \ \vdash\ f(u) = f(v)$
    \item All 1138 declarations in 5 seconds; a ``Formula'' fold in the index.
  \end{itemize}
  \textbf{Remaining problems}
  \begin{itemize}\small
    \item Constants without notation print as $\mathrm{isGroup}(g)$.
    \item Conjunctions nest with explicit brackets: $A \land (B \land (C \land \dots))$.
    \item Claude's conventions not yet confirmed; long formulas only scroll.
  \end{itemize}
\end{frame}

\begin{frame}{Task 10: GodementSyntax --- why}
  \begin{itemize}
    \item Informath's grammar covers the statements that come \emph{out} of Dedukti.
    \item Godement \emph{writes} differently. Counted in the 599 formalized sentences
          (UD trees, regex + dependency patterns):
  \end{itemize}
  \begin{center}\footnotesize
  \begin{tabular}{lrl}
    \toprule
    \textbf{Construction} & \textbf{sentences} & \textbf{in the grammar?} \\
    \midrule
    ``$x$ is said to be $A$ if \dots'' & 53 & \bad{no} \\
    ``$x$ is called the $N$ of $y$'' & 65 & \bad{no} \\
    ``for $x$ to be $A$ it is necessary and sufficient that'' & 24 & \bad{no} \\
    ``the following conditions are equivalent: a) \dots{} b) \dots'' & 33 & \bad{no} \\
    ``there exists a unique $x$ such that'' & 22 & \bad{no} \\
    ``subgroup of $G$ which contains $B$'' (relative clause) & 105 & \bad{no} \\
    ``if $P$ then $Q$'', ``for all $x$'', ``there exists'' & 180+ & \ok{yes} \\
    \bottomrule
  \end{tabular}
  \end{center}
\end{frame}

\begin{frame}{Task 10: GodementSyntax --- prompt and interaction}
  \prompt{Let us investigate the linguistic structures used in definitions and
    theorem statements in Godement. [\dots] Can you make a list of the structures
    of statements in Godement by using the Informath categories Prop, Kind, Exp,
    Adj, etc? [\dots] Do not try to run GF in parsing mode yet, because it can be slow.}
  \prompt{Excellent! Can you suggest GF abstract syntax functions,
    linearizations using RGL, and mappings to MathCore}
  \prompt{You can add GodementSyntax to UserExtensions and a test case [\dots]
    and a Makefile entry for godement\_syntax that produces a pdf and opens it.}
  \begin{itemize}\small
    \item Parsing real sentences revealed what linearizing did not:
          \texttt{\$...\$} is indexed before parsing, a hidden argument crashes the
          parser, a comma too many blocks a parse.
  \end{itemize}
\end{frame}

\begin{frame}{Task 10: GodementSyntax --- outcome}
  \commit{865eae0} (26 Sept): about 50 functions; semantics in Haskell;
  a flag \texttt{-translate-core}; \texttt{make godement\_syntax} $\to$ PDF.
  \begin{itemize}\small
    \item \emph{For an element $x$ of $G$ to be reflexible, it is necessary and
          sufficient that $x$ has a left reflection.}\\
          $\Rightarrow$ for all elements $x$ of $G$, ($x$ is reflexible, if and only if
          there is a left reflection of $x$).
    \item \emph{There exists a unique element $x$ of $G$ such that $x$ is invertible.}\\
          $\Rightarrow$ there exists $x$ \dots{} and for all $y$ \dots, if $y$ is
          invertible, then $y = x$.
    \item 38 of 38 test sentences parse; no regressions on the old tests.
  \end{itemize}
  \textbf{Remaining problems}
  \begin{itemize}\small
    \item Ambiguity: several readings per sentence; ugly fresh variables \texttt{\_g0}.
    \item English only; ``$x$ is a $K$'' has no Dedukti translation when $K$ is a type.
    \item Not yet run on Godement's actual 10\,413 sentences.
  \end{itemize}
\end{frame}

\begin{frame}{Task 11: GodementExamples --- why}
  \begin{itemize}
    \item Symbol tables speak Godement only as far as \texttt{Examples.gf} lets them.
    \item The formalization notes (written with Claude) listed the gaps:
      \begin{itemize}\small
        \item no three-place noun: ``a homomorphism of $G$ into $H$'' became
              \bad{``$f$ is a homomorphism''};
        \item no ``into'', ``onto'', ``generated by'';
        \item no negation: ``$x$ is not an element of $y$'';
        \item ``zero'' and ``one'' but no ``two''.
      \end{itemize}
    \item 216 table entries showed fewer arguments than their constants had.
  \end{itemize}
\end{frame}

\begin{frame}{Task 11: GodementExamples --- prompt and interaction}
  \prompt{We noticed that some syntactic forms that should be usable in .dkgf
    entries are missing from the grammar. [\dots] I would like to start with a new
    module GodementExamples [\dots] You can try to rebuild the dkgf files by using
    these new functions where necessary. If this means that you need some new
    categories, add them and their application functions also to GodementExamples.}
  \begin{itemize}\small
    \item New categories need more than grammar: the Dedukti $\leftrightarrow$ MathCore
          conversions in Haskell had to learn them.
    \item Verified by reading chapters back with old and new tables and diffing.
    \item \emph{``go on''}: a second round of rewordings, and the whole index rebuilt.
  \end{itemize}
\end{frame}

\begin{frame}{Task 11: GodementExamples --- outcome}
  \commit{35184d0} (26 Sept): 3 new categories, 5 example forms, 12 prepositions.
  \begin{center}\small
  \begin{tabular}{ll}
    \toprule
    \textbf{before} & \textbf{after} \\
    \midrule
    \texttt{"\#1 is a homomorphism"} & \texttt{"\#1 is a homomorphism of \#2 into \#3"} \\
    (not worded) & \texttt{"the subgroup of \#1 generated by \#2"} \\
    \texttt{"the class of \#3 under \#2"} & \texttt{"the class of \#3 modulo \#2 in \#1"} \\
    \texttt{\$\#1 \textbackslash notin \#2\$} only & \texttt{"\#1 is not an element of \#2"} \\
    \bottomrule
  \end{tabular}
  \end{center}
  \textbf{Remaining problems}
  \begin{itemize}\small
    \item ``the set of all subsets of $X$ with $n$ elements'' still does not parse.
    \item Negated definitions read as tautologies: ``$x \notin y$ if $x \notin y$''.
    \item The next morning: English readings of §§ 10--24 had never been possible;
          21 declarations are too big to word in 30 seconds.
  \end{itemize}
\end{frame}

% =====================================================================
\section{Observations and the future}

\begin{frame}{What we observed}
  \begin{columns}[T]
    \begin{column}{0.48\textwidth}
      \ok{\textbf{Worked well}}
      \begin{itemize}\small
        \item Bulk lexical work with parallel subagents (2000 terms in minutes).
        \item Tireless crash-driven iteration (Czech RGL).
        \item Self-audits \emph{when asked}: the Czech paradigm check.
        \item Grounding claims in measurements and diffs.
        \item Code in the style of the surrounding code.
      \end{itemize}
    \end{column}
    \begin{column}{0.48\textwidth}
      \bad{\textbf{Needed a human}}
      \begin{itemize}\small
        \item Design decisions: unique macro names, symbolic preference,
              ``just time out and move on''.
        \item \texttt{linking ... OK} $\neq$ correct.
        \item Careless moments: files swept into a commit, an installed binary
              overwritten mid-comparison.
        \item Native-speaker judgement: none of the new words is reviewed.
      \end{itemize}
    \end{column}
  \end{columns}
  \vspace{0.5em}
  \centering\small
  The division of labour: I set direction and judge; Claude explores, builds, measures.
\end{frame}

\begin{frame}{Speculation: what would be good to do next}
  \begin{itemize}\small
    \item \textbf{Review}: native checking of the Finnish, Czech, Polish and French
          words; gender agreement of ``such that'' in Slavic languages.
    \item \textbf{Tests}: a regression suite and CI --- today regressions are found
          by diffing outputs by hand.
    \item \textbf{Scale}: faster or incremental search for wordings, so that
          big statements (determinants of order 3) get English too.
    \item \textbf{Proofs}: repair truncated proofs, English theorem headings,
          proper bound-variable names.
    \item \textbf{Coverage}: run GodementSyntax on all 10\,413 sentences and measure
          how much of the book parses; port it to French, German, Swedish.
    \item \textbf{Autoformalization}: parse Godement's own text into Dedukti and
          type-check it against the hand-made formalization.
    \item \textbf{Hybrid}: an LLM proposes symbol-table entries and grammar rules,
          the grammar and the type checker verify them.
  \end{itemize}
\end{frame}

\begin{frame}[standout]
  Thank you!

  \vspace{1em}
  \normalsize
  \texttt{github.com/GrammaticalFramework/informath}
\end{frame}

\end{document}
